\chapter{Curves over Finite Fields}
\begin{center}\itshape
How many solutions has $x^3+y^3=1$ modulo $p$?\\
About $p$, give or take $2\sqrt p$.\\
It took from Gauss to Weil to see why it is never more.
\end{center}

In Article 358 of the \emph{Disquisitiones Arithmeticae} (1801), Gauss counted the solutions of a cubic congruence, in effect the points of the
curve $x^3+y^3=1$ modulo a prime $p\equiv1\pmod3$, and found an exact formula in terms of the representation $4p=L^2+27M^2$. Thirteen years later, on
9~July 1814, he wrote the last of the 146 entries of his mathematical diary: if $a+bi$ is a Gaussian prime with $a-1+bi$ divisible by $2+2i$, the number of
solutions of $1=x^2+y^2+x^2y^2$ modulo $a+bi$ is $(a-1)^2+b^2$. In both cases the number of points is $p$ plus an error that is at most about $2\sqrt p$.

A century later, in his 1921 thesis, Emil Artin attached a zeta function to curves $y^2=f(x)$ over $\mathbb F_p$, imitating Dedekind's, and checked
on some forty examples that its zeros satisfy the analogue of the Riemann hypothesis. In 1932 Harold Davenport and Louis Mordell were bounding character
sums, which count points on curves; Artin told Hasse that the best possible bound, over all finite extensions of the field, is exactly the Riemann hypothesis for
these zeta functions, and early in 1933 Hasse proved it for elliptic curves. The general case was the work of Andr\'e Weil. In 1940, held in the military prison
of Rouen for failing to report for service, he worked so well that he wrote to his wife that he might have to arrange to spend two or three months locked
up every year; the note he sent from there to the \emph{Comptes Rendus} announced the Riemann hypothesis for all curves over finite fields, and the proof
appeared in 1948. In 1949, recalling Gauss's cubic congruences, Weil conjectured the same for varieties of every dimension. Deligne proved this in 1974.

Weil's proof used the geometry of correspondences on $C\times C$ and of the Jacobian of the curve. In 1969 Sergei Stepanov found an elementary method for $y^2=f(x)$, using an
auxiliary polynomial with many zeros, the trick of transcendence theory. Enrico Bombieri turned it in 1973 into a two-page proof of the upper bound for all
curves, using nothing but the Riemann--Roch theorem. This chapter proves the rationality and the functional equation of the zeta function of a smooth plane
curve with the Riemann--Roch theorem of Chapter~11, Bombieri's upper bound, and the full Riemann hypothesis for the curves $y^m=f(x)$, which include
elliptic curves (Hasse's theorem) and Fermat curves.

\begin{reminder}[from Chapter 11]
$k$ is an algebraically closed field; here $k=\overline{\mathbb F}_q$. $C$ is a smooth plane curve $F=0$ of degree $n$ and genus $g=\frac{(n-1)(n-2)}2$, with function
field $k(C)$, valuations $\operatorname{ord}_P$, divisors, $L(D)$ and $\ell(D)=\dim_kL(D)$. Principal divisors have degree $0$, and a function without poles is
constant (Proposition~\ref{prop:divbasic}); $\ell(D+P)\le\ell(D)+1$ (Proposition~\ref{prop:ellbasic}); Riemann--Roch, $\ell(D)-\ell(W-D)=\deg D+1-g$ with
$W=(n-3)E$, $E=\operatorname{div}(L)$ for a line $L$ (Theorem~\ref{thm:rr}); in particular $\ell(D)=\deg D+1-g$ when $\deg D>2g-2$. If $\operatorname{div}(G)\ge\operatorname{div}(H)$
for forms $G,H$ prime to $F$, then $G=AF+BH$ with $B$ a form (Corollary~\ref{cor:noethersmooth}). The finite field $\mathbb F_{q^r}$ is the set of roots of
$x^{q^r}-x$ in $k$, and $x\mapsto x^q$ is an automorphism of $k$ fixing exactly $\mathbb F_q$ (Volume~2).
\end{reminder}

\section{Frobenius and rational divisors}

Throughout, $q$ is a power of a prime $p$, $k=\overline{\mathbb F}_q$, and $C$ is a smooth plane curve $F=0$ of degree $n$ whose equation $F$ has coefficients in
$\mathbb F_q$; we say that $C$ is \emph{defined over $\mathbb F_q$}. The \emph{Frobenius} of $\mathbb P^2$ is
\[
\varphi(a:b:c)=(a^q:b^q:c^q).
\]
It is a bijection (with inverse the $q$-th root, $k$ being perfect), and $F(\varphi P)=F(P)^q$ because the coefficients of $F$ are fixed by $x\mapsto x^q$, so
$\varphi$ maps $C$ onto $C$. A point is fixed by $\varphi^r$ iff it has a representative with coordinates in $\mathbb F_{q^r}$ (normalise its first non-zero
coordinate to $1$). Let
\[
N_r=|C(\mathbb F_{q^r})|=\#\{P\in C:\varphi^rP=P\}.
\]
Every point of $C$ has algebraic coordinates, so it lies in some $C(\mathbb F_{q^r})$: the orbits of $\varphi$ are finite. An orbit is a \emph{closed point} $\mathfrak P$,
and its size is its \emph{degree} $\deg\mathfrak P$. A point $P$ lies in $C(\mathbb F_{q^r})$ iff the size of its orbit divides $r$, so, with $B_d$ the number of closed
points of degree $d$,
\begin{equation}\label{eq:NrBd}
N_r=\sum_{d\mid r}dB_d.
\end{equation}
A divisor $D=\sum n_PP$ is \emph{rational} if $\varphi_*D=D$, where $\varphi_*D=\sum n_P\,\varphi(P)$: its coefficients are constant on orbits, so the rational
divisors are the integer combinations of closed points (each $\mathfrak P$ read as the sum of its points).

\subsection*{Frobenius on functions}

For a form $G$, let $G^\sigma$ be the form obtained by raising every coefficient of $G$ to the power $q$; then $G^\sigma(\varphi P)=G(P)^q$, and $F^\sigma=F$.
For $z=G/H\in k(C)$ set $z^\sigma=G^\sigma/H^\sigma$. This is well defined, since $F\mid GH'-G'H$ implies $F=F^\sigma\mid(GH'-G'H)^\sigma$, and
$z\mapsto z^\sigma$ is a bijective ring map of $k(C)$ with $(cz)^\sigma=c^qz^\sigma$ for constants $c$. The functions \emph{defined over $\mathbb F_q$} are
\[
k(C)_q=\{G/H:\ G,H\text{ forms with coefficients in }\mathbb F_q\},
\]
a subfield of $k(C)$ fixed by $\sigma$, whose constants are $k\cap k(C)_q=\mathbb F_q$ (a constant $c$ with $c^q=c$ is in $\mathbb F_q$).

\begin{lemma}\label{lem:frobdiv}
For $z\in k(C)^\times$, $\operatorname{div}(z^\sigma)=\varphi_*\operatorname{div}(z)$. If $z$ is regular at $P$, then $z^\sigma$ is regular at $\varphi P$ and
$z^\sigma(\varphi P)=z(P)^q$. In particular, if $z\in k(C)_q$ and $P\in C(\mathbb F_q)$ is not a pole of $z$, then $z(P)\in\mathbb F_q$.
\end{lemma}
\begin{proof}
It suffices to treat a form $G$ prime to $F$. Let $P$ be a point, say with $Z\ne0$ (then $\varphi P$ too), and work with dehomogenised polynomials. The ring
automorphism $h\mapsto h^\sigma$ of $k[x,y]$ satisfies $h^\sigma(\varphi P)=h(P)^q$, so it maps the local ring $\mathcal O_P(\mathbb P^2)$ onto
$\mathcal O_{\varphi P}(\mathbb P^2)$ and the ideal $(F,G)$ onto $(F,G^\sigma)$. It induces an additive bijection
$\mathcal O_P(\mathbb P^2)/(F,G)\to\mathcal O_{\varphi P}(\mathbb P^2)/(F,G^\sigma)$ with $\sigma(cv)=c^q\sigma(v)$, which maps a $k$-basis to a $k$-basis (if
$\sum c_i\sigma(e_i)=0$, apply $\sigma^{-1}$). So $I(\varphi P,F\cap G^\sigma)=I(P,F\cap G)$, which is the first claim. If $z=G/H$ with $H(P)\ne0$, then
$H^\sigma(\varphi P)=H(P)^q\ne0$ and $z^\sigma(\varphi P)=G(P)^q/H(P)^q$. For the last claim, $z^\sigma=z$ and $\varphi P=P$ give $z(P)=z(P)^q$.
\end{proof}

\subsection*{Descent}

The space $L(D)$ is a vector space over $k$. When $D$ is rational, we want the functions of $L(D)$ that are defined over $\mathbb F_q$, and to know that there
are as many as there should be. Everything rests on a piece of linear algebra.

\begin{lemma}\label{lem:descentlin}
Let $W\subseteq k^M$ be a subspace with $\sigma(W)\subseteq W$, where $\sigma$ raises each coordinate to the power $q$. Then $W$ has a basis of vectors with
coordinates in $\mathbb F_q$.
\end{lemma}
\begin{proof}
Let $s=\dim W$. Some $s$ coordinates $I\subseteq\{1,\dots,M\}$ are independent on $W$, that is, the projection $\pi_I:W\to k^I$ is an isomorphism (take $s$
independent columns of a matrix whose rows form a basis of $W$). For $i\in I$ let $w_i\in W$ be the vector with $\pi_I(w_i)=e_i$, the $i$-th unit vector.
Then $\sigma(w_i)\in W$ and $\pi_I(\sigma(w_i))=\sigma(e_i)=e_i$, so $\sigma(w_i)=w_i$: the coordinates of $w_i$ satisfy $c^q=c$, so they lie in
$\mathbb F_q$. The $w_i$ form a basis of $W$.
\end{proof}

\begin{lemma}\label{lem:normform}
For every rational divisor $D$ there is a form $H_0$ with coefficients in $\mathbb F_q$, not divisible by $F$, with $\operatorname{div}(H_0)\ge D$.
\end{lemma}
\begin{proof}
It suffices to find, for each closed point $\mathfrak P$ of degree $d$, such a form $H_{\mathfrak P}$ vanishing at all points of $\mathfrak P$; a product of powers of
these then dominates $D$. Let $P\in\mathfrak P$; its coordinates can be taken in $\mathbb F_{q^d}$. Let $M$ be a line through $P$ with coefficients in $\mathbb F_{q^d}$,
other than $C$ itself if $C$ is a line (for $P=(a:b:1)$, one of $X-aZ$, $Y-bZ$ works; similarly in the other charts). Put
\[
H_{\mathfrak P}=M\,M^\sigma M^{\sigma^2}\cdots M^{\sigma^{d-1}}.
\]
Since $M^{\sigma^d}=M$, $\sigma$ permutes the factors, so $H_{\mathfrak P}^\sigma=H_{\mathfrak P}$ and its coefficients lie in $\mathbb F_q$. The factor $M^{\sigma^j}$
vanishes at $\varphi^jP$. If $F\mid H_{\mathfrak P}$, then $F$, being irreducible, would be a scalar multiple of some $M^{\sigma^j}$, so $n=1$ and $C=M^{\sigma^j}$;
applying $\sigma^{-j}$ (which fixes $F$) gives $C=M$, excluded.
\end{proof}

\begin{proposition}[descent]\label{prop:descent}
Let $D$ be a rational divisor.
\begin{enumerate}[label=(\alph*)]
\item The space $L_q(D)=L(D)\cap k(C)_q$ over $\mathbb F_q$ has dimension $\ell(D)$, and every $\mathbb F_q$-basis of it is a $k$-basis of $L(D)$.
\item If $z\in k(C)$ satisfies $z^\sigma=z$, then $z\in k(C)_q$.
\item If $D'$ is a rational divisor with $D'\equiv D$, then $D'=D+\operatorname{div}(w)$ for some $w\in k(C)_q^\times$.
\end{enumerate}
\end{proposition}
\begin{proof}
(a) \emph{Elements of $k(C)_q$ that are independent over $\mathbb F_q$ are independent over $k$.} Otherwise take a relation $\sum c_iz_i=0$ with the fewest
non-zero $c_i$, and scale it so that some $c_{i_0}=1$. Applying $\sigma$ gives $\sum c_i^qz_i=0$, and subtracting, $\sum(c_i^q-c_i)z_i=0$ has fewer non-zero
coefficients (the $i_0$-th is $0$). So all $c_i^q=c_i$, the $c_i$ lie in $\mathbb F_q$, and the $z_i$ are dependent over $\mathbb F_q$. Hence
$\dim_{\mathbb F_q}L_q(D)\le\ell(D)$.

\emph{Enough functions.} Take $H_0$ as in Lemma~\ref{lem:normform}, of degree $d$, and $D_0=\operatorname{div}(H_0)-D\ge0$, a rational divisor. Let $R_d$ be the
forms of degree $d$ and
\[
V=\{G\in R_d:\ G/H_0\in L(D)\}=\{G:\ F\mid G\text{ or }\operatorname{div}(G)\ge D_0\}.
\]
Every $z\in L(D)$ is of the form $G/H_0$ with $G\in V$: write $z=G_1/H_1$; then
$\operatorname{div}(G_1H_0)=(\operatorname{div}(z)+D)+D_0+\operatorname{div}(H_1)\ge\operatorname{div}(H_1)$, so $G_1H_0=AF+GH_1$ with $G\in R_d$
(Corollary~\ref{cor:noethersmooth}), and $z=G/H_0$ in $k(C)$. So $G\mapsto G/H_0$ maps $V$ onto $L(D)$, with kernel $U=F\cdot R_{d-n}$, and
$\ell(D)=\dim V-\dim U$. Both $V$ and $U$ are stable under $\sigma$: if $\operatorname{div}(G)\ge D_0$ then
$\operatorname{div}(G^\sigma)=\varphi_*\operatorname{div}(G)\ge\varphi_*D_0=D_0$ by Lemma~\ref{lem:frobdiv}. Let $V_q$ and $U_q$ be the forms of $V$ and $U$ with
coefficients in $\mathbb F_q$. By Lemma~\ref{lem:descentlin}, $V_q$ contains a $k$-basis of $V$, so $\dim_{\mathbb F_q}V_q\ge\dim V$; and vectors with entries in
$\mathbb F_q$ that are independent over $\mathbb F_q$ are independent over $k$ (the rank of a matrix is computed by Gaussian elimination inside $\mathbb F_q$), so
$\dim_{\mathbb F_q}V_q=\dim V$, and likewise $\dim_{\mathbb F_q}U_q=\dim U$. The map $V_q\to L_q(D)$, $G\mapsto G/H_0$, has kernel $V_q\cap U=U_q$, so
$\dim_{\mathbb F_q}L_q(D)\ge\dim V-\dim U=\ell(D)$. With the first step, $\dim_{\mathbb F_q}L_q(D)=\ell(D)$, and an $\mathbb F_q$-basis of $L_q(D)$ is a family of
$\ell(D)$ elements independent over $k$: a $k$-basis of $L(D)$.

(b) By Lemma~\ref{lem:frobdiv}, $\operatorname{div}(z)$ is rational, so the divisor $D$ of poles of $z$ is rational and $z\in L(D)$. Write $z=\sum c_iz_i$ in an
$\mathbb F_q$-basis of $L_q(D)$, which is a $k$-basis of $L(D)$ by (a). Applying $\sigma$, $z=\sum c_i^qz_i$, so $c_i^q=c_i$, $c_i\in\mathbb F_q$, and
$z\in k(C)_q$.

(c) Write $D'-D=\operatorname{div}(z)$. Then $\operatorname{div}(z^\sigma)=\varphi_*(D'-D)=\operatorname{div}(z)$, so $z^\sigma/z$ has no zeros or poles: it is a
constant $\lambda\ne0$ (Proposition~\ref{prop:divbasic}). As $k$ is algebraically closed, there is $\mu$ with $\mu^{q-1}=\lambda^{-1}$, and
$w=\mu z$ satisfies $w^\sigma=\mu^q\lambda z=w$. By (b), $w\in k(C)_q$, and $\operatorname{div}(w)=\operatorname{div}(z)$.
\end{proof}

Part (c) is a disguised form of Hilbert's Theorem~90: the ``cocycle'' $\lambda$ is a coboundary because every element of $k^\times$ is a $(q-1)$-th power.

\section{The zeta function}

\begin{definition}
Let $A_d$ be the number of effective rational divisors of degree $d$. The \emph{zeta function} of $C$ over $\mathbb F_q$ is the power series
\[
Z(T)=\sum_{d\ge0}A_dT^d.
\]
\end{definition}

$A_d$ is finite, since an effective rational divisor of degree $d$ is a sum of closed points of degree $\le d$, and there are at most $N_j$ closed points of degree
$j$. Setting $T=q^{-s}$, $\zeta_C(s)=Z(q^{-s})=\sum_D N(D)^{-s}$ with $N(D)=q^{\deg D}$ is the exact analogue of the Dedekind zeta function of Chapter~8, a sum over
the ``ideals'' $D$ of their norms.

\begin{proposition}[Euler product]\label{prop:euler}
As formal power series,
\[
Z(T)=\prod_{\mathfrak P}\frac1{1-T^{\deg\mathfrak P}}=\exp\Bigl(\sum_{r\ge1}\frac{N_r}rT^r\Bigr).
\]
\end{proposition}
\begin{proof}
An effective rational divisor is $\sum_{\mathfrak P}n_{\mathfrak P}\mathfrak P$ with $n_{\mathfrak P}\ge0$, uniquely; expanding the product gives each divisor of
degree $d$ once as a term $T^d$ (only finitely many closed points have degree $\le d$, so each coefficient is a finite sum). Taking the logarithmic derivative,
\[
T\frac{Z'(T)}{Z(T)}=\sum_{\mathfrak P}\frac{\deg\mathfrak P\,T^{\deg\mathfrak P}}{1-T^{\deg\mathfrak P}}=\sum_{\mathfrak P}\sum_{j\ge1}\deg\mathfrak P\,T^{j\deg\mathfrak P}
=\sum_{r\ge1}\Bigl(\sum_{d\mid r}dB_d\Bigr)T^r,
\]
which is $\sum_rN_rT^r$ by \eqref{eq:NrBd}. Both sides of the claimed identity have constant term $1$ and the same logarithmic derivative.
\end{proof}

For a line, $C=\mathbb P^1$, $N_r=q^r+1$ and
\[
Z(T)=\exp\Bigl(\sum_r\frac{(qT)^r}r+\sum_r\frac{T^r}r\Bigr)=\frac1{(1-T)(1-qT)}.
\]
Exercise~\ref{ex:zetaP1} recovers it from the definition. The theorem below says that every curve is a modification of this one by a polynomial.

\begin{lemma}\label{lem:classcount}
Two rational divisors are \emph{equivalent} if they differ by $\operatorname{div}(w)$, $w\in k(C)_q^\times$; by Proposition~\ref{prop:descent}(c), this is linear
equivalence. Let $D$ be rational. The number of effective rational divisors equivalent to $D$ is $\dfrac{q^{\ell(D)}-1}{q-1}$.
\end{lemma}
\begin{proof}
They are the divisors $D+\operatorname{div}(w)$ with $w\in L_q(D)$, $w\ne0$, and $w,w'$ give the same divisor iff $w'/w$ is a constant in $k(C)_q$, that is, an
element of $\mathbb F_q^\times$. By Proposition~\ref{prop:descent}(a), $L_q(D)$ has $q^{\ell(D)}$ elements.
\end{proof}

\begin{theorem}[F.~K.~Schmidt; rationality and functional equation]\label{thm:rationality}
Let $C$ be a smooth plane curve of genus $g$ defined over $\mathbb F_q$.
\begin{enumerate}[label=(\alph*)]
\item There are rational divisors of degree $1$, and the number $h$ of classes of rational divisors of degree $0$ is finite.
\item $Z(T)=\dfrac{P(T)}{(1-T)(1-qT)}$, where $P$ is a polynomial of degree $2g$ with integer coefficients, $P(0)=1$ and $P(1)=h$.
\item $Z\bigl(\frac1{qT}\bigr)=q^{1-g}T^{2-2g}Z(T)$; equivalently $P(T)=q^gT^{2g}P\bigl(\frac1{qT}\bigr)$.
\item Writing $P(T)=\prod_{i=1}^{2g}(1-\alpha_iT)$ with $\alpha_i\in\C$, the map $\alpha\mapsto q/\alpha$ permutes the $\alpha_i$ (with multiplicities), and
for all $r\ge1$,
\[
N_r=q^r+1-\sum_{i=1}^{2g}\alpha_i^r.
\]
\end{enumerate}
\end{theorem}
\begin{proof}
\emph{Finiteness.} The degrees of rational divisors form a subgroup $\delta\Z$ of $\Z$, $\delta\ge1$; fix a rational $D_1$ of degree $\delta$. Subtracting
$jD_1$ is a bijection from the classes of degree $j\delta$ to the classes of degree $0$. When $j\delta\ge g$, every class of degree $j\delta$ contains an effective
divisor (Riemann's inequality gives $\ell(D)\ge1$, and Lemma~\ref{lem:classcount}), so there are at most $A_{j\delta}$ such classes. Hence $h\le A_{j\delta}$ is finite.

\emph{The series.} $A_d=0$ unless $\delta\mid d$, and grouping the effective divisors of degree $d$ by class (Lemma~\ref{lem:classcount}),
\begin{equation}\label{eq:Adclasses}
A_d=\sum_{[D],\ \deg D=d}\frac{q^{\ell(D)}-1}{q-1}.
\end{equation}
For $d>2g-2$, $\ell(D)=d+1-g$ for every $D$ of degree $d$ (Riemann--Roch), so $A_d=h\,\frac{q^{d+1-g}-1}{q-1}$ when moreover $\delta\mid d$. So
\[
Z(T)=\sum_{d\le2g-2}A_dT^d+\frac h{q-1}\sum_{\substack{d>2g-2\\\delta\mid d}}\bigl(q^{d+1-g}-1\bigr)T^d.
\]
The second sum is a combination of geometric series in $T^\delta$ and $(qT)^\delta$, so $Z$ is a rational function. Near $T=1$, the series in $(qT)^\delta$ is
holomorphic (its poles are at $|T|=1/q$), while $\sum_{\delta\mid d,\,d>2g-2}T^d=T^{d_0}/(1-T^\delta)$ has a simple pole at $T=1$, where $d_0$ is the least such $d$. As $h\ge1$, $Z$ has a \emph{simple pole at $T=1$}.

\emph{$\delta=1$.} Let $Z_\delta$ be the zeta function of $C$ over $\mathbb F_{q^\delta}$; its numbers of points are $N_{\delta s}$. For any $r$, with $\zeta$ running over
the $r$-th roots of unity,
\[
\prod_{\zeta}Z(\zeta T)=\exp\Bigl(\sum_{n\ge1}\frac{N_n}nT^n\sum_\zeta\zeta^n\Bigr)=\exp\Bigl(\sum_{s\ge1}\frac{N_{rs}}{s}T^{rs}\Bigr)=Z_r(T^r),
\]
since $\sum_\zeta\zeta^n$ is $r$ if $r\mid n$ and $0$ otherwise. Take $r=\delta$: as $Z(T)$ is a series in $T^\delta$, $Z(\zeta T)=Z(T)$, and $Z_\delta(T^\delta)=Z(T)^\delta$.
The right side has a pole of order $\delta$ at $T=1$. The left side has a simple pole there, because $Z_\delta$ has a simple pole at $1$ (what was just proved, over
$\mathbb F_{q^\delta}$) and $T\mapsto T^\delta$ has a simple zero of $T^\delta-1$ at $T=1$. So $\delta=1$.

\emph{The polynomial.} If $g=0$, then $A_d=h\frac{q^{d+1}-1}{q-1}$ for all $d\ge0$, so $Z(T)=\frac h{q-1}\bigl(\frac q{1-qT}-\frac1{1-T}\bigr)=\frac h{(1-T)(1-qT)}$,
and $Z(0)=1$ gives $h=1$, $P=1$; (c) and (d) are immediate. Let $g\ge1$. With $\delta=1$, split \eqref{eq:Adclasses} as $Z(T)=\frac1{q-1}F_1(T)+G(T)$ with
\[
F_1(T)=\sum_{[D],\ 0\le\deg D\le2g-2}q^{\ell(D)}T^{\deg D},
\qq
G(T)=\frac h{q-1}\Bigl(\frac{q^gT^{2g-1}}{1-qT}-\frac1{1-T}\Bigr);
\]
indeed $G$ collects the terms $-\frac h{q-1}T^d$ for $0\le d\le2g-2$ and $\frac h{q-1}(q^{d+1-g}-1)T^d$ for $d\ge2g-1$. Then $P(T)=(1-T)(1-qT)Z(T)$ is a polynomial
of degree $\le2g$, with $P(0)=A_0=1$, and integer coefficients since $Z$ has them. At $T=1$, $(1-T)(1-qT)F_1(T)$ vanishes and $(1-T)(1-qT)G(T)\to\frac h{q-1}(q-1)=h$:
so $P(1)=h$.

\emph{Functional equation.} Choose the line $L$ defining $W=(n-3)\operatorname{div}(L)$ with coefficients in $\mathbb F_q$ (one of $X,Y,Z$ is not $C$); then $W$ is
rational, of degree $2g-2$. By Riemann--Roch, $q^{\ell(D)}=q^{\deg D+1-g}q^{\ell(W-D)}$, and $[D]\mapsto[W-D]$ permutes the classes of degree between $0$ and
$2g-2$, sending degree $d$ to $2g-2-d$. So
\[
F_1\Bigl(\frac1{qT}\Bigr)=\sum_{[D]}q^{1-g}q^{\ell(W-D)}T^{-\deg D}=q^{1-g}T^{2-2g}\sum_{[D]}q^{\ell(W-D)}T^{\deg(W-D)}=q^{1-g}T^{2-2g}F_1(T).
\]
A direct computation gives the same for $G$:
\[
G\Bigl(\frac1{qT}\Bigr)=\frac h{q-1}\Bigl(\frac{qT}{1-qT}-\frac{q^{1-g}T^{2-2g}}{1-T}\Bigr)=q^{1-g}T^{2-2g}G(T).
\]
This is (c) for $Z$. Since $(1-\frac1{qT})(1-\frac1T)=\frac{(1-qT)(1-T)}{qT^2}$, it becomes $P(T)=q^gT^{2g}P(\frac1{qT})$. Comparing the coefficients of $T^{2g}$,
the leading coefficient of $P$ is $q^gP(0)=q^g\ne0$, so $\deg P=2g$.

(d) If $P(1/\alpha)=0$, then $P(\alpha/q)=q^g(\alpha/q)^{2g}P(1/\alpha)=0$, so $q/\alpha$ is an $\alpha_j$; since (c) is an identity of polynomials, multiplicities
match. Finally, taking logarithms in $Z=P/((1-T)(1-qT))$ and comparing with Proposition~\ref{prop:euler},
\[
\sum_r\frac{N_r}rT^r=-\sum_r\frac{\sum_i\alpha_i^r}rT^r+\sum_r\frac{T^r}r+\sum_r\frac{q^rT^r}r.
\]
\end{proof}

\begin{remark}
Since $\prod\alpha_i=q^g$ and the $\alpha_i$ come in pairs $\alpha,q/\alpha$, the statement $|\alpha_i|=\sqrt q$ for all $i$ is equivalent to: all zeros of
$\zeta_C(s)=Z(q^{-s})$ lie on the line $\operatorname{Re}s=\frac12$ (a zero is $q^{s}=\alpha_i$). This is the \emph{Riemann hypothesis for $C$}. It implies the
\emph{Hasse--Weil bound}
\[
|N_r-q^r-1|\le2g\,q^{r/2}.
\]
For an elliptic curve $E$, $P(T)=1-aT+qT^2$ with $a=q+1-N_1$, and $h=P(1)=N_1$: the classes of degree $0$ are the points of $E(\mathbb F_q)$, as
Corollary~\ref{cor:pic} predicts.
\end{remark}

Over $\mathbb F_q$, the numbers $N_1,\dots,N_g$ determine $P$, hence all $N_r$: the coefficients of $T^j$ for $j\le g$ come from $N_1,\dots,N_j$ by Newton's
identities, and (c) gives the others. A curve over a finite field knows its point counts over all extensions from finitely many of them.

\section{Counting with characters}

Before any general theory, let us count points on one curve exactly, as Gauss did.

\begin{proposition}[Gauss]\label{prop:gausscubic}
Let $p\ne3$ be a prime and $N$ the number of points of $X^3+Y^3=Z^3$ over $\mathbb F_p$. If $p\equiv2\pmod3$, $N=p+1$. If $p\equiv1\pmod3$, then
\[
N=p+1+L,
\]
where $L$ is the unique integer with $L\equiv1\pmod3$ and $4p=L^2+27M^2$ for some integer $M$. In particular $|N-p-1|<2\sqrt p$.
\end{proposition}
\begin{proof}
The curve is smooth (its partial derivatives $3X^2,3Y^2,-3Z^2$ vanish together only at $0$). If $p\equiv2\pmod3$, cubing is a bijection of $\mathbb F_p$
(its kernel on $\mathbb F_p^\times$ is trivial), so $y^3=1-x^3$ has one solution $y$ for each $x$, and at infinity $X^3=-Y^3$ gives the one point $(1:-1:0)$.

Let $p\equiv1\pmod3$, and let $\chi$ be a character of order $3$ of $\mathbb F_p^\times$, extended by $\chi(0)=0$, and $\chi^0$ the trivial character with
$\chi^0(0)=1$. The number of $x$ with $x^3=a$ is $1+\chi(a)+\overline\chi(a)$: for $a=0$ it is $1$, and for $a\ne0$ it is $3$ if $a$ is a cube and $0$ otherwise.
So the affine count is
\[
\sum_{a+b=1}\Bigl(\sum_{i=0}^2\chi^i(a)\Bigr)\Bigl(\sum_{j=0}^2\chi^j(b)\Bigr)=\sum_{i,j}J(\chi^i,\chi^j),\qquad J(\psi,\psi')=\sum_{a+b=1}\psi(a)\psi'(b).
\]
Here $J(\chi^0,\chi^0)=p$, $J(\chi^0,\chi^j)=\sum_b\chi^j(b)=0$ for $j\ne0$, and $J(\chi,\overline\chi)=\sum_{a\ne0,1}\chi\bigl(\frac a{1-a}\bigr)=-\chi(-1)=-1$,
because $a\mapsto\frac a{1-a}$ maps $\mathbb F_p\setminus\{0,1\}$ onto $\mathbb F_p\setminus\{0,-1\}$ and $\chi(-1)=1$. So the affine count is
$p-2+J(\chi,\chi)+\overline{J(\chi,\chi)}$. At infinity, $(X/Y)^3=-1$ has three solutions, since $-1$ is a cube and $\mathbb F_p$ has three cube roots of $1$.
Hence $N=p+1+2\operatorname{Re}J(\chi,\chi)$.

Now use Chapter~10, Lemma~\ref{lem:jacobi}, with $\chi=\chi_\pi$ for a primary prime $\pi$ of norm $p$ in $D=\Z[\omega]$ (any cubic character is $\chi_\pi$ or its
conjugate, which give the same $N$): $J(\chi,\chi)=\pi=a+b\omega$ with $a\equiv2$, $b\equiv0\pmod3$. Then $2\operatorname{Re}\pi=2a-b=:L\equiv1\pmod3$ and
$4p=4(a^2-ab+b^2)=(2a-b)^2+3b^2=L^2+27M^2$ with $b=3M$. For uniqueness, if $4p=L'^2+27M'^2$ with $L'\equiv1\pmod 3$, then $L'\equiv M'\pmod2$ and
$\beta=\frac{L'+3M'}2+3M'\omega\in D$ has $2\operatorname{Re}\beta=L'$ and norm $\frac{L'^2+27M'^2}4=p$. So $\beta$ is an associate of $\pi$ or $\overline\pi$
($D$ is a UFD), and it is primary ($\frac{L'+3M'}2\equiv2L'\equiv2\pmod3$); by Lemma~\ref{lem:primary}, $\beta\in\{\pi,\overline\pi\}$ and $L'=L$. Finally
$L^2<4p$ since $M\ne0$ ($4p$ is not a square).
\end{proof}

For $p=7$: $28=1+27$, $L=1$, $N=9$. For $p=13$: $52=25+27$, $L=-5$, $N=9$. For $p=31$: $124=16+108$, $L=4$, $N=36$. The bound $|N-p-1|\le2\sqrt p$ fell out
of the arithmetic of $\Z[\omega]$ for free: it is the Riemann hypothesis for this curve, of genus $1$, proved a century before anyone formulated it. Exercise~\ref{ex:jacobicount}
does the same for $y^2=x^3+1$.

\begin{think}
For the curve $y^2=x^3-x+1$, the number of points over $\mathbb F_p$ is $p+1+\sum_x\bigl(\frac{x^3-x+1}p\bigr)$. If the $p$ Legendre symbols were independent fair
coin tosses $\pm1$, how large would the sum typically be, and how often would it exceed $2\sqrt p$? Compare with Figure~\ref{fig:hasse}: in $667$ primes the
dots never leave the parabola.
\end{think}

\begin{figure}[ht]
\centering\fig{ch13-hasse}
\caption{$a_p=p+1-N_p$ for $y^2=x^3-x+1$ and the primes $3\le p<5000$, $p\ne23$. All dots lie between the parabolas $\pm2\sqrt p$ (Corollary~\ref{cor:hasse}); the largest
ratio $|a_p|/2\sqrt p$ is $0.977$.}
\label{fig:hasse}
\end{figure}

\section{Bombieri's upper bound}

The idea of Stepanov's method is the idea of transcendence proofs: to show that a set $S$ is small, build a non-zero function that vanishes to high order at every
point of $S$ and has few poles. Its degree, zeros minus poles, is $0$, so $|S|$ times the order of vanishing is at most the number of poles. The difficulty is to
make a function vanish to high order at $q$ points without paying for it: the Frobenius does that for free, because $v^q$ is a $q$-th power.

\begin{theorem}[Stepanov, Bombieri]\label{thm:bombieri}
Let $C$ be a smooth plane curve of genus $g$ defined over $\mathbb F_q$, where $q=q_0^2$ is a square and $q>(g+1)^4$. Then
\[
N_1\le q+(2g+1)\sqrt q.
\]
\end{theorem}
\begin{proof}
Assume $N_1\ge1$ and fix $P_0\in C(\mathbb F_q)$. Put $m=q_0-1$ and $n=q_0+2g$.

\emph{A basis with distinct pole orders.} For $j=0,1,\dots,m$, the spaces $L_q(jP_0)$ increase, and their dimensions $\ell(jP_0)$ grow by at most $1$ at each step
(Proposition~\ref{prop:ellbasic}, Proposition~\ref{prop:descent}(a)). For each $j$ where the dimension grows, choose $u\in L_q(jP_0)\setminus L_q((j-1)P_0)$
(take $u=1$ for $j=0$): it has a pole of order exactly $j$ at $P_0$. This gives $u_1,\dots,u_\rho\in k(C)_q$ with distinct pole orders $s_i\in[0,m]$ at $P_0$, no
other poles, and $\rho=\ell(mP_0)\ge m+1-g=q_0-g$ by Riemann's inequality.

\emph{The auxiliary function.} Consider the $\mathbb F_q$-linear map
\[
\Phi:\ L_q(nP_0)^\rho\to L_q\bigl((q_0m+n)P_0\bigr),\qquad(v_1,\dots,v_\rho)\mapsto\sum_{i=1}^\rho u_i^{q_0}v_i.
\]
As $n>2g-2$ and $q_0m+n=q+2g$, Riemann--Roch and Proposition~\ref{prop:descent}(a) give the dimensions $n+1-g=q_0+g+1$ and $q+g+1$. Now
\[
\rho(q_0+g+1)-(q+g+1)\ge(q_0-g)(q_0+g+1)-q-g-1=q_0-(g+1)^2>0,
\]
because $q_0>(g+1)^2$. So $\Phi$ has a non-zero kernel element $(v_1,\dots,v_\rho)$, and we set
\[
A=\sum_{i=1}^\rho u_i^{q_0}v_i^{q}=\Bigl(\sum_{i=1}^\rho u_iv_i^{q_0}\Bigr)^{q_0},
\]
the second form because raising to the power $q_0$, a power of $p$, is additive.

\emph{$A\ne0$.} At $P_0$, the term $u_iv_i^{q_0}$ with $v_i\ne0$ has a pole of order $s_i+q_0t_i$, where $t_i\ge0$ is the pole order of $v_i$. Since $0\le s_i<q_0$
and the $s_i$ are distinct, these orders are distinct for distinct $i$; so the term of largest order cannot be cancelled, and $\sum u_iv_i^{q_0}\ne0$.

\emph{Zeros.} Let $P\in C(\mathbb F_q)$, $P\ne P_0$. All $u_i,v_i$ are regular at $P$, and $v_i(P)\in\mathbb F_q$ by Lemma~\ref{lem:frobdiv}, so $v_i(P)^q=v_i(P)$
and
\[
A(P)=\sum u_i(P)^{q_0}v_i(P)^q=\sum u_i(P)^{q_0}v_i(P)=\Phi(v)(P)=0.
\]
As $A$ is a $q_0$-th power of a function regular at $P$, $\operatorname{ord}_P(A)\ge q_0$.

\emph{Counting.} $A$ has poles only at $P_0$, of order at most $q_0m+qn$. Since $\deg\operatorname{div}(A)=0$,
\[
q_0(N_1-1)\le q_0m+qn,\qquad N_1\le1+m+q_0n=q+(2g+1)q_0.
\]
\end{proof}

The bound is applied to $C$ over $\mathbb F_{q^r}$ for every even $r$ with $q^r>(g+1)^4$, where it gives $N_r\le q^r+(2g+1)q^{r/2}$: the $\alpha_i$ cannot conspire to make
$-\sum\alpha_i^r$ much larger than $q^{r/2}$. That is only half of what is needed.

\section{The Riemann hypothesis}

\begin{lemma}\label{lem:boundsRH}
Let $\alpha_1,\dots,\alpha_s\in\C$ and $S_r=\sum_i\alpha_i^r$. If there are $K$ and $r_1$ with $|S_r|\le Kq^{r/2}$ for all even $r\ge r_1$, then $|\alpha_i|\le\sqrt q$
for all $i$.
\end{lemma}
\begin{proof}
Let $\beta_i=\alpha_i^2$ and group equal values: $\sum_{t\ge1}S_{2t}T^t=\sum_\beta m_\beta\frac{\beta T}{1-\beta T}$, the sum over the distinct non-zero $\beta$, with
multiplicities $m_\beta\ge1$. The left side converges for $|T|<1/q$, since $|S_{2t}T^t|\le K(q|T|)^t$ for $2t\ge r_1$. The right side is a rational function with a
pole at each $1/\beta$ (the residues $-m_\beta/\beta$ do not cancel), and its Taylor series at $0$ converges exactly up to the nearest pole. So $|1/\beta|\ge1/q$ for
every $\beta$: $|\alpha_i|^2\le q$.
\end{proof}

With Theorem~\ref{thm:rationality}(d), an upper bound $|\alpha_i|\le\sqrt q$ for all $i$ is the whole Riemann hypothesis: the pairing $\alpha\mapsto q/\alpha$ then gives
$|\alpha_i|\ge\sqrt q$ too. But Lemma~\ref{lem:boundsRH} needs a \emph{two-sided} bound on $S_r=q^r+1-N_r$, and Theorem~\ref{thm:bombieri} bounds $N_r$ only from above.
This is not a technicality: $\alpha=\frac{7+\sqrt{13}}2$, $\overline\alpha=\frac{7-\sqrt{13}}2$ satisfy $\alpha\overline\alpha=9$, and $q^r+1-\alpha^r-\overline\alpha^r\le9^r+1$
for all $r$, yet $|\alpha|>3$. A curve with too few points is as bad as one with too many. For a family of curves, the missing half comes from counting
all the twists at once.

\begin{proposition}[twists]\label{prop:twists}
Let $m\ge2$ with $p\nmid m$, and $f\in\mathbb F_Q[x]$ without repeated roots, of degree $m$ or $m+1$. For $c\in\mathbb F_Q^\times$, the projective closure $C_c$ of
$y^m=cf(x)$ is a smooth plane curve of degree $\max(m,\deg f)$, defined over $\mathbb F_Q$. Let $e=\gcd(m,Q-1)$ and $c_1,\dots,c_e$ representatives of the cosets
of the subgroup $H=\{y^m:y\in\mathbb F_Q^\times\}$ of index $e$. Then
\[
\sum_{j=1}^e|C_{c_j}(\mathbb F_Q)|=e(Q+1).
\]
\end{proposition}
\begin{proof}
\emph{Smoothness.} At an affine singular point, $my^{m-1}=0$ and $cf'(x)=0$, so $y=0$, $f(x)=0=f'(x)$: impossible. If $\deg f=m+1$, the equation is
$Y^mZ=cf^*(X,Z)$, which meets $Z=0$ only at $(0:1:0)$; there the derivative in $Z$ is $1$, since every term of $f^*$ has degree $m+1\ge2$ in $X,Z$. If
$\deg f=m$, the equation is $Y^m=cf^*(X,Z)$; at infinity $Y^m=ca_mX^m$, the points are $(1:\zeta:0)$ with $\zeta^m=ca_m\ne0$, and the derivative in $Y$ is
$m\zeta^{m-1}\ne0$.

\emph{Counting.} The group $\mathbb F_Q^\times$ is cyclic of order $Q-1$, so $y\mapsto y^m$ has kernel of order $e$ and image $H$ of index $e$. Hence for
$b\in\mathbb F_Q$ the equation $y^m=b$ has $1$ solution if $b=0$, $e$ if $b\in H$, $0$ otherwise. For fixed $x$ with $f(x)\ne0$, exactly one of the $c_jf(x)$
lies in $H$, and if $f(x)=0$ each curve has the single point $(x,0)$; so the $e$ curves together have $e$ affine points above each $x$, $eQ$ in all. At infinity: if $\deg f=m+1$, each curve has one point; if $\deg f=m$,
the curve $C_{c_j}$ has $e$ points if $c_ja_m\in H$ and none otherwise, and exactly one $j$ has $c_ja_m\in H$. In both cases the points at infinity add up to $e$.
\end{proof}

\begin{theorem}[Riemann hypothesis for $y^m=f(x)$]\label{thm:rhfamily}
Let $C$ be the curve $y^m=f(x)$ over $\mathbb F_q$, with $m,f$ as in Proposition~\ref{prop:twists}, of genus $g$. Then $|\alpha_i|=\sqrt q$ for $i=1,\dots,2g$.
\end{theorem}
\begin{proof}
Let $r$ be even with $Q=q^r>(g+1)^4$; $Q$ is a square. Apply Proposition~\ref{prop:twists} over $\mathbb F_Q$ with $c_1=1$, so that $C_{c_1}=C$ and
$|C(\mathbb F_Q)|=N_r$. Every $C_{c_j}$ is a smooth plane curve over $\mathbb F_Q$ of the same degree, hence of genus $g$, so Theorem~\ref{thm:bombieri} gives
$|C_{c_j}(\mathbb F_Q)|\le Q+(2g+1)\sqrt Q$. For $j=1$ this is the upper bound, and the twist identity gives the lower one:
\[
\begin{aligned}
N_r=e(Q+1)-\sum_{j\ge2}|C_{c_j}(\mathbb F_Q)|&\ge e(Q+1)-(e-1)\bigl(Q+(2g+1)\sqrt Q\bigr)\\
&\ge Q+1-(e-1)(2g+1)\sqrt Q.
\end{aligned}
\]
So $|S_r|=|q^r+1-N_r|\le m(2g+1)q^{r/2}$ for all such $r$. By Lemma~\ref{lem:boundsRH}, $|\alpha_i|\le\sqrt q$, and by the pairing of
Theorem~\ref{thm:rationality}(d), $|\alpha_i|=\sqrt q$.
\end{proof}

\begin{corollary}[Hasse, 1933]\label{cor:hasse}
For $C$ as in Theorem~\ref{thm:rhfamily}, $|N_r-q^r-1|\le2g\,q^{r/2}$ for all $r$. In particular, if $q$ is odd and $f$ is a cubic without repeated roots,
the elliptic curve $y^2=f(x)$ has
\[
\bigl|\,|E(\mathbb F_q)|-q-1\,\bigr|\le2\sqrt q.
\]
\end{corollary}
\begin{proof}
$|\sum_i\alpha_i^r|\le2g\,q^{r/2}$. For the elliptic curve, $m=2$, $\deg f=3=m+1$ and $g=1$.
\end{proof}

The Fermat curves $x^m+y^m=1$ ($f=1-x^m$) are covered too, and so the curve of Proposition~\ref{prop:gausscubic}, whose bound we already knew.

\begin{beauty}
The proof of the Riemann hypothesis for these curves uses no complex analysis at all, only Riemann--Roch, the Frobenius and a dimension count; the complex numbers
appear only at the end, to say where the roots of a polynomial with integer coefficients lie. And the Frobenius, the map $v\mapsto v^q$, plays two opposite roles:
it is the symmetry whose fixed points are being counted, and it is the tool that makes the auxiliary function vanish to high order at all of them at once.
\end{beauty}

\begin{remark}[Weil's theorem]
The Riemann hypothesis holds for every smooth projective curve over $\mathbb F_q$ (Weil, announced in 1940, proved in 1948 with the intersection theory of
correspondences on $C\times C$). Theorem~\ref{thm:bombieri} holds for all of them, with the same proof, once Riemann--Roch is known for all smooth curves. What is
missing here for a general plane curve is the lower bound. Bombieri obtains it from a Galois covering $C'\to\mathbb P^1$ through which $C$ factors: the
twists of $C'$ by the elements of the Galois group play the role of the curves $C_{c_j}$, and the count on $\mathbb P^1$, exactly $Q+1$, plays the role of
$e(Q+1)$. The curve $C'$ is in general not a plane curve, which is why we stop at the family $y^m=f(x)$. We use Weil's theorem in general only in remarks.
\end{remark}

\begin{remark}[the Weil conjectures]
For a smooth projective variety $X$ of dimension $d$ over $\mathbb F_q$, Weil conjectured in 1949 that $Z_X(T)=\exp(\sum|X(\mathbb F_{q^r})|T^r/r)$ is rational,
satisfies a functional equation, has the shape $\prod_{j=0}^{2d}P_j(T)^{(-1)^{j+1}}$ with the zeros of $P_j$ of absolute value $q^{-j/2}$, and that the degrees of the
$P_j$ are the Betti numbers of the corresponding complex variety. Dwork proved rationality in 1960 by $p$-adic analysis; Grothendieck built \'etale cohomology to
explain the rest, and Deligne proved the Riemann hypothesis in 1974. One consequence is Ramanujan's conjecture $|\tau(p)|\le2p^{11/2}$ for the coefficients of the
discriminant modular form (Chapter~14).
\end{remark}

\begin{remark}[Sato--Tate]
By Corollary~\ref{cor:hasse}, for an elliptic curve over $\Q$ and a good prime $p$, $a_p=p+1-N_p=2\sqrt p\cos\theta_p$ for an angle $\theta_p\in[0,\pi]$. How are
the angles distributed? Figure~\ref{fig:satotate} shows $\cos\theta_p$ for $y^2=x^3-x+1$ and the $78\,496$ primes $p<10^6$ other than $23$: not uniform, but along a
semicircle. Sato and Tate conjectured in the 1960s that for curves without complex multiplication, the $\theta_p$ are equidistributed for the measure
$\frac2\pi\sin^2\theta\,\mathrm d\theta$; this was proved by Clozel, Harris, Shepherd-Barron and Taylor (2008) for curves with non-integral $j$-invariant, and by
Barnet-Lamb, Geraghty, Harris and Taylor (2011) for all of them. For curves with complex multiplication, half the $a_p$ vanish (Exercise~\ref{ex:supersing}).
\end{remark}

\begin{figure}[ht]
\centering\fig{ch13-satotate}
\caption{Histogram of $a_p/2\sqrt p$ for $y^2=x^3-x+1$, $p<10^6$ ($20$ bins), and the Sato--Tate density $\frac2\pi\sqrt{1-x^2}$.}
\label{fig:satotate}
\end{figure}

\section{Exercises}

\begin{exercise}[$\star$]\label{ex:zetaP1}
For $C=\mathbb P^1$ over $\mathbb F_q$, show that $A_d=\frac{q^{d+1}-1}{q-1}$ by counting directly (an effective divisor of degree $d$ is the zero set of a binary form of
degree $d$ over $\mathbb F_q$, up to a scalar), and recover $Z(T)=\frac1{(1-T)(1-qT)}$. Deduce from Theorem~\ref{thm:rationality} that every smooth plane conic defined
over $\mathbb F_q$ has exactly $q^r+1$ points over $\mathbb F_{q^r}$.
\end{exercise}

\begin{exercise}[$\star$]\label{ex:ellNr}
For an elliptic curve with $N_1=q+1-a$, show that $N_r=q^r+1-\alpha^r-\overline\alpha^r$, where $\alpha,\overline\alpha$ are the roots of $x^2-ax+q$. Count the points of
$y^2=x^3-x+1$ over $\mathbb F_5$ by hand, and deduce $N_2$ and $N_3$ without counting.
\end{exercise}

\begin{exercise}[$\star\star$]\label{ex:jacobicount}
Let $p\equiv1\pmod3$, $\rho$ the Legendre symbol and $\chi$ a cubic character modulo $p$. Show that $y^2=x^3+1$ has $p+1+2\operatorname{Re}J(\rho,\chi)$ points over
$\mathbb F_p$, and that $|J(\rho,\chi)|=\sqrt p$ (show $g(\rho)g(\chi)=J(\rho,\chi)g(\rho\chi)$ and use $|g(\psi)|=\sqrt p$, Volume~5, Theorem~3.5). Deduce Hasse's bound for this
curve.
\end{exercise}

\begin{exercise}[$\star\star$]\label{ex:moment}
For $a,b\in\mathbb F_p$ ($p$ odd) let $s(a,b)=\sum_{x\in\mathbb F_p}\bigl(\frac{x^3+ax+b}p\bigr)$. Show that
\[
\sum_{a,b\in\mathbb F_p}s(a,b)^2=p^2(p-1).
\]
Interpret: on average over all curves, $a_p^2$ is about $p$. What does the Sato--Tate measure predict for the average of $(2\sqrt p\cos\theta)^2$?
\end{exercise}

\begin{exercise}[$\star\star$]\label{ex:supersingp2}
Let $p\equiv2\pmod3$, $p>2$. Using Exercise~\ref{ex:supersing}, find the zeta function of $y^2=x^3+1$ over $\mathbb F_p$, and show that over $\mathbb F_{p^2}$ it has
$(p+1)^2$ points. Which values of $r$ give $N_r=p^r+1$?
\end{exercise}

\begin{exercise}[$\star\star$]\label{ex:maximal}
Let $q$ be a power of $p$ and $C$ the Fermat curve $X^{q+1}+Y^{q+1}=Z^{q+1}$ over $\mathbb F_{q^2}$. Show that it has genus $\frac{q(q-1)}2$ and $q^3+1$ points over
$\mathbb F_{q^2}$ (use the norm map $x\mapsto x^{q+1}$ from $\mathbb F_{q^2}^\times$ onto $\mathbb F_q^\times$). Deduce from Theorem~\ref{thm:rhfamily} that $C$ is
\emph{maximal}: $P(T)=(1+qT)^{2g}$. (This curve is isomorphic to the Hermitian curve of Exercise~\ref{ex:hermitian}.)
\end{exercise}

\begin{exercise}[$\star\star\star$]\label{ex:klein}
The Klein quartic is $X^3Y+Y^3Z+Z^3X=0$, over $\mathbb F_2$. Show that it is smooth, of genus $3$, that $N_1=3$ and $N_2=5$, and, given $N_3=24$ (check by counting points over $\mathbb F_8$), that $P(T)=1+5T^3+8T^6$. Deduce $N_r$ for all $r$, and show that over $\mathbb F_8$ the curve has $q+1+g\lfloor2\sqrt q\rfloor$ points, the most allowed by
Serre's refinement of the Hasse--Weil bound.
\end{exercise}

\section{Hints}

\begin{hint}{ex:zetaP1}
Binary forms of degree $d$: $q^{d+1}-1$ non-zero ones. For the conic, $g=0$, so $P=1$.
\end{hint}
\begin{hint}{ex:ellNr}
$P(T)=1-aT+qT^2=(1-\alpha T)(1-\overline\alpha T)$. Over $\mathbb F_5$ the squares are $0,1,4$.
\end{hint}
\begin{hint}{ex:jacobicount}
Count $y^2=u$, $x^3=-v$ with $u+v=1$, as in Proposition~\ref{prop:gausscubic}. For the Gauss sums, expand $g(\rho)g(\chi)$ and group by $s+t$.
\end{hint}
\begin{hint}{ex:moment}
Expand the square. For $x\ne x'$, $(a,b)\mapsto(x^3+ax+b,x'^3+ax'+b)$ is a bijection of $\mathbb F_p^2$.
\end{hint}
\begin{hint}{ex:supersingp2}
$N_1=p+1$ gives $a=0$, so $\alpha=i\sqrt p$.
\end{hint}
\begin{hint}{ex:maximal}
The equation is $y^{q+1}=1-x^{q+1}$. If $\sum_{i=1}^{2g}\alpha_i=-2gq$ and $|\alpha_i|=q$, every $\alpha_i$ is $-q$.
\end{hint}
\begin{hint}{ex:klein}
In characteristic $2$, a singular point has $X^2Y=Z^3$, $X^3=Y^2Z$, $Y^3=Z^2X$; show $X^7=Y^7=Z^7$, then evaluate $F$. Over $\mathbb F_4$, $t^3=1$ for $t\ne0$.
For $P$: Newton's identities and the functional equation.
\end{hint}

\section{Solutions}

\begin{solution}{ex:zetaP1}
\idea Divisors on $\mathbb P^1$ are zero sets of binary forms; then compare with the formula of Proposition~\ref{prop:euler}.
\writeup A closed point of $\mathbb P^1$ of degree $d$ is the zero set of an irreducible binary form $h(X,Z)$ of degree $d$ over $\mathbb F_q$, up to a scalar:
$Z$ for the point at infinity, $Z^d\pi(X/Z)$ with $\pi$ monic irreducible otherwise (the orbit of a root of $\pi$). So an effective rational divisor of degree $d$ is a
non-zero binary form of degree $d$ up to $\mathbb F_q^\times$, by unique factorisation: $A_d=\frac{q^{d+1}-1}{q-1}$. Then
\[
Z(T)=\frac1{q-1}\sum_d(q\cdot q^dT^d-T^d)=\frac1{q-1}\Bigl(\frac q{1-qT}-\frac1{1-T}\Bigr)=\frac1{(1-T)(1-qT)},
\]
as found from $N_r=q^r+1$. A smooth conic has $n=2$, $g=0$, so $P$ has degree $0$ and $P(0)=1$: $P=1$, and Theorem~\ref{thm:rationality}(d) gives
$N_r=q^r+1$. In particular every smooth conic over $\mathbb F_q$ has a rational point, and then (projecting from it) it is a copy of $\mathbb P^1$.
\end{solution}

\begin{solution}{ex:ellNr}
\idea With $g=1$, $P$ has two coefficients determined by $N_1$ and the functional equation.
\writeup $P(T)=1+c_1T+c_2T^2$ with $c_2=q$ by Theorem~\ref{thm:rationality}(c) (leading coefficient $q^g$), and $N_1=q+1-(\alpha+\overline\alpha)=q+1+c_1$, so
$c_1=-a$ and $P=1-aT+qT^2$, whose inverse roots are the roots of $x^2-ax+q$. Theorem~\ref{thm:rationality}(d) gives $N_r$.

For $f(x)=x^3-x+1$ over $\mathbb F_5$: $f(0)=1$, $f(1)=1$, $f(2)=7\equiv2$, $f(3)=25\equiv0$, $f(4)=61\equiv1$. The squares modulo $5$ are $0,1,4$, so the number of $y$
is $2,2,0,1,2$: $7$ affine points, $N_1=8$, $a=-2$. Then $\alpha+\overline\alpha=-2$, $\alpha\overline\alpha=5$,
\[
\alpha^2+\overline\alpha^2=4-10=-6,\qquad\alpha^3+\overline\alpha^3=(-2)^3-3\cdot5\cdot(-2)=22,
\]
so $N_2=26+6=32$ and $N_3=126-22=104$ (numerical evaluation confirms both).
\end{solution}

\begin{solution}{ex:jacobicount}
\idea Each count $\#\{y:y^2=u\}$, $\#\{x:x^3=v\}$ is a sum of characters, and the cross terms are Jacobi sums.
\writeup The affine points are the pairs with $y^2=u$, $x^3=-v$, $u+v=1$; replacing $x$ by $-x$, $\#\{x:x^3=-v\}=\#\{x:x^3=v\}=1+\chi(v)+\overline\chi(v)$, and
$\#\{y:y^2=u\}=1+\rho(u)$. Expanding,
\[
\sum_{u+v=1}(1+\rho(u))(1+\chi(v)+\overline\chi(v))=p+J(\rho,\chi)+J(\rho,\overline\chi),
\]
since $\sum_u\rho(u)=\sum_v\chi(v)=0$. Also $J(\rho,\overline\chi)=\overline{J(\rho,\chi)}$, and the point at infinity adds $1$: $N=p+1+2\operatorname{Re}J(\rho,\chi)$.

For non-trivial $\psi,\psi'$ with $\psi\psi'$ non-trivial, $g(\psi)g(\psi')=\sum_u\zeta^u\sum_{s+t=u}\psi(s)\psi'(t)$. For $u\ne0$, substituting $s=us'$, $t=ut'$, the
inner sum is $\psi\psi'(u)J(\psi,\psi')$; for $u=0$ it is $\psi'(-1)\sum_s\psi\psi'(s)=0$. So $g(\psi)g(\psi')=J(\psi,\psi')g(\psi\psi')$, and taking absolute values
(Volume~5, Theorem~3.5), $|J(\psi,\psi')|=\sqrt p\sqrt p/\sqrt p=\sqrt p$. With $\psi=\rho$, $\psi'=\chi$ ($\rho\chi$ has order $6$), $|N-p-1|\le2\sqrt p$.
\end{solution}

\begin{solution}{ex:moment}
\idea Squaring a character sum and summing over the family turns it into a count of pairs $(x,x')$, and most pairs cancel.
\writeup $\sum_{a,b}s(a,b)^2=\sum_{x,x'}\sum_{a,b}\rho(x^3+ax+b)\rho(x'^3+ax'+b)$. For $x=x'$ the inner sum counts the $(a,b)$ with $x^3+ax+b\ne0$: $p^2-p$. For
$x\ne x'$, the map $(a,b)\mapsto(x^3+ax+b,x'^3+ax'+b)$ is affine with determinant $x-x'\ne0$, a bijection of $\mathbb F_p^2$, so the inner sum is
$(\sum_u\rho(u))^2=0$. Total: $p(p^2-p)=p^2(p-1)$. For the $p^2$ pairs, the mean of $a_p^2=s^2$ is $p-1$. Under Sato--Tate,
$\int_0^\pi4p\cos^2\theta\cdot\frac2\pi\sin^2\theta\,\mathrm d\theta=\frac{2p}\pi\int_0^\pi\sin^22\theta\,\mathrm d\theta=p$: the same.
\end{solution}

\begin{solution}{ex:supersingp2}
\idea $N_1$ gives $P$, and $P$ gives everything.
\writeup By Exercise~\ref{ex:supersing}, $N_1=p+1$, so $a=0$, $P(T)=1+pT^2$, $Z(T)=\frac{1+pT^2}{(1-T)(1-pT)}$, and $\alpha=i\sqrt p$, $\overline\alpha=-i\sqrt p$.
Then $\alpha^r+\overline\alpha^r=(i\sqrt p)^r(1+(-1)^r)$: it is $0$ for odd $r$ and $2(-p)^{r/2}$ for even $r$. So $N_r=p^r+1$ for odd $r$, and for $r=2$,
$N_2=p^2+1+2p=(p+1)^2$; for $r\equiv0\pmod4$, $N_r=p^r+1-2p^{r/2}=(p^{r/2}-1)^2$, the minimum allowed.
\end{solution}

\begin{solution}{ex:maximal}
\idea Count with the norm map; a point count at the Weil bound forces all $\alpha_i$ to be equal.
\writeup The curve has degree $q+1$, so $g=\frac{q(q-1)}2$, and it is $y^{q+1}=f(x)$ with $f=1-x^{q+1}$, $m=q+1$ prime to $p$, $f$ without repeated roots
($f'=-x^q$ vanishes only at $0$, and $f(0)=1$). The map $N(x)=x^{q+1}$ sends $\mathbb F_{q^2}^\times$ onto $\mathbb F_q^\times$ (its image has order
$\frac{q^2-1}{q+1}=q-1$ and lies in $\mathbb F_q$, since $N(x)^q=N(x)$), each fibre having $q+1$ elements, and $N(0)=0$. The affine points are the pairs with
$N(x)=u$, $N(y)=v$, $u+v=1$ in $\mathbb F_q$: the pairs $(0,1)$ and $(1,0)$ give $q+1$ points each, and the $q-2$ others $(q+1)^2$ each, in all
$(q+1)\bigl(2+(q-2)(q+1)\bigr)=q^3-q$. At infinity, $N(X/Y)=-1$ has $q+1$ solutions. So $N_1=q^3+1=q^2+1+2gq$.

By Theorem~\ref{thm:rationality}(d), $\sum\alpha_i=q^2+1-N_1=-2gq$, and $|\alpha_i|=q$ by Theorem~\ref{thm:rhfamily} (over $\mathbb F_{q^2}$, $\sqrt{q^2}=q$). Then
$\sum\operatorname{Re}\alpha_i=-2gq$ with each $\operatorname{Re}\alpha_i\ge-q$ forces $\alpha_i=-q$ for all $i$: $P(T)=(1+qT)^{2g}$, and $N_r=q^{2r}+1-2g(-q)^r$.
\end{solution}

\begin{solution}{ex:klein}
\idea Smoothness by direct elimination; then three point counts and the functional equation pin down the six coefficients of $P$.
\writeup \emph{Smoothness.} In characteristic $2$, $F_X=X^2Y+Z^3$, $F_Y=X^3+Y^2Z$, $F_Z=Y^3+Z^2X$. If $X=0$ at a singular point, then $Z=0$ and $Y=0$; so
by cyclic symmetry all coordinates are non-zero, and we may take $Z=1$. Then $Y=X^{-2}$ from $F_X$, and $F_Y$ gives $X^3=X^{-4}$, $X^7=1$. Now
$F=X^3Y+Y^3+X=X+X^{-6}+X=X^{-6}=X\ne0$, so the point is not on the curve. The genus is $3$.

\emph{Counts.} Over $\mathbb F_2$, the points with a zero coordinate are $(1:0:0),(0:1:0),(0:0:1)$ (if $Z=0$ then $X^3Y=0$), and $(1:1:1)$ gives $F=1$: $N_1=3$. Over
$\mathbb F_4$, the points with all coordinates non-zero and $Z=1$ satisfy $X^3=Y^3=1$, so $F=Y+1+X$: $X+Y=1$ with $X,Y\ne0$ gives $(\omega,\omega^2)$ and
$(\omega^2,\omega)$, so $N_2=5$.

\emph{Zeta.} $s_r=\sum\alpha_i^r=q^r+1-N_r$: $s_1=0$, $s_2=0$, $s_3=9-24=-15$. With $P(T)=1-e_1T+e_2T^2-e_3T^3+\dots$, Newton's identities give $e_1=s_1=0$,
$e_2=\frac{s_1^2-s_2}2=0$, $e_3=\frac{s_1^3-3s_1s_2+2s_3}6=-5$. The functional equation $c_{6-j}=q^{3-j}c_j$ gives $c_4=2c_2=0$, $c_5=4c_1=0$, $c_6=8$:
$P(T)=1+5T^3+8T^6$. So $P(T)=Q(T^3)$ with $Q(U)=1+5U+8U^2$, the $\alpha_i$ are the cube roots of the inverse roots $\beta,\overline\beta$ of $Q$, and
$N_r=2^r+1$ when $3\nmid r$ (the sum of the $r$-th powers of the three cube roots of $\beta$ vanishes), while $N_{3s}=8^s+1-3(\beta^s+\overline\beta^s)$, with
$\beta+\overline\beta=-5$, $\beta\overline\beta=8$. For $s=1$, $N_3=9+15=24$, and $\lfloor2\sqrt8\rfloor=5$, so $N_3=8+1+3\cdot5$.
\end{solution}

\section*{Sources and further reading}

\begin{itemize}
\item C.~F.~Gauss, \emph{Disquisitiones Arithmeticae}, 1801, Article~358; the diary entry of 9~July 1814 is no.~146 of his \emph{Tagebuch}. See also K.~Ireland and
M.~Rosen, \emph{A Classical Introduction to Modern Number Theory}, 2nd ed., Springer GTM~84, 1990, Chapters~8 and~11, for the Jacobi sum counts.
\item P.~Roquette, ``The Riemann hypothesis in characteristic $p$, its origin and development. Part~1'', \emph{Mitt.\ Math.\ Ges.\ Hamburg} 21 (2002); Part~2, 22 (2004).
Artin's thesis, Davenport and Mordell, Artin's remark to Hasse in 1932, and Hasse's proof of 1933.
\item A.~Weil, ``Sur les fonctions alg\'ebriques \`a corps de constantes fini'', \emph{C.\ R.\ Acad.\ Sci.\ Paris} 210 (1940), 592--594; ``Numbers of solutions of
equations in finite fields'', \emph{Bull.\ AMS} 55 (1949), 497--508. Weil's letter from Rouen prison to his wife is quoted in the MacTutor History of Mathematics
archive, ``Andr\'e Weil writes from Rouen prison''.
\item S.~A.~Stepanov, ``On the number of points of a hyperelliptic curve over a finite prime field'', \emph{Izv.\ Akad.\ Nauk SSSR Ser.\ Mat.} 33 (1969), 1171--1181.
\item E.~Bombieri, ``Counting points on curves over finite fields (d'apr\`es S.~A.~Stepanov)'', S\'eminaire Bourbaki, expos\'e 430 (1972/73), Springer LNM~383, 1974.
The proof of Theorem~\ref{thm:bombieri}, and the lower bound by Galois coverings.
\item H.~Stichtenoth, \emph{Algebraic Function Fields and Codes}, 2nd ed., Springer GTM~254, 2009, Chapter~5. Rationality, functional equation and the Riemann
hypothesis for all curves, in the language of function fields; Serre's bound.
\item W.~M.~Schmidt, \emph{Equations over Finite Fields: An Elementary Approach}, Springer LNM~536, 1976. Stepanov's method in full generality.
\item P.~Deligne, ``La conjecture de Weil.~I'', \emph{Publ.\ Math.\ IH\'ES} 43 (1974), 273--307.
\item L.~Clozel, M.~Harris and R.~Taylor, ``Automorphy for some $\ell$-adic lifts of automorphic mod~$\ell$ Galois representations'', \emph{Publ.\ Math.\ IH\'ES} 108 (2008);
T.~Barnet-Lamb, D.~Geraghty, M.~Harris and R.~Taylor, ``A family of Calabi--Yau varieties and potential automorphy~II'', \emph{Publ.\ RIMS} 47 (2011), 29--98.
The Sato--Tate conjecture.
\end{itemize}
